% ---------------------------------------------------------------------------
% Internal-model ablation: LDNet vs linear unsteady aerodynamics.
% Draft for chapter3.tex, to follow subsec:mpc_results.
% Numbers from light/ablation/ (see NOTES.md). Table body generated by
% collect.py into table_body.tex -> ablation_table_body.tex on Overleaf.
%
% Style checked against light/latex/AGENTS.md:
%   - "plant" is forbidden: the controlled system is the "reduced system"
%   - chord is c_ref, gust amplitude r_g = W_0/U_inf, effort weight R / R*
%   - Test 2 values written once at definition (already in 4.2), label after
%   - no \paragraph, no \emph, no explanatory parentheses
%   - citations: Theodorsen, Jones1940, hoblit1988gust already in the .bib
% ---------------------------------------------------------------------------

\subsection{Ablation of the internal model}
\label{subsec:ablation}

The results of \cref{subsec:mpc_results} rest on the premise that the
alleviation requires an internal model able to reproduce a nonlinear,
viscous response. We test that premise by replacing the retained LDNet
inside the same loop with a linear unsteady aerodynamic model, at equal
preview and equal cost function, and we measure what the substitution costs.

The substitute is the classical finite-state model of the theories cited in
\cref{sec:introduction}: circulatory lift from the effective downwash
filtered through the Wagner indicial function in the two-pole approximation
of Jones~\cite{Jones1940}, gust lift filtered through the K\"ussner
function in a two-pole approximation~\cite{hoblit1988gust}, and the
non-circulatory contribution of the apparent mass~\cite{Theodorsen}. Its
four lag states are integrated alongside the structural state, and it
exposes the same interface as the LDNet, so that the Model Predictive
Controller, the flap grid, the rate and saturation limits, the preview and
the structural integrator are untouched: the two arms differ only in the
model the controller uses.

We fit its coefficients by least squares to the same campaign that trained
the LDNet, rather than taking them from thin-airfoil theory. A linear model
carrying textbook gains would be penalised for a gain error instead of for
the physics it omits, and the comparison would not measure what it claims
to. Fitting gives it the best coefficients available within its own
structure. Where the campaign does not identify a coefficient we retain the
theoretical value: the airfoil-motion downwash contributes only
\SI{2.8}{\percent} of the standard deviation of $C_L$, so the lift-curve
slope is weakly determined, and pinning it to $2\pi$ costs \num{0.0015} in
fit NRMSE.

\Cref{tab:ablation} reports both arms over the CS-25.341 grid, each tuned
independently with the selection rule of \cref{subsec:mpc_results} so that
neither is judged at the other's operating point. The LDNet arm reproduces
\cref{tab:cs25_grid} to within \SI{0.05}{pp}. The linear arm does not
alleviate the gust: averaged over the grid it increases the peak lift
excursion, and on several cells it violates the stability check.

That failure is not a matter of tuning. Sweeping $R$ over six decades on
Test~2 leaves no value at which the linear controller both acts and helps:
below $R=10^{-2}$ it violates the $\dot\alpha$ bound with a pitch ratio
above $4$, and above it the effort penalty drives the commanded deflection
to zero, so that the best value passing the stability check, $R=1$, returns
exactly \SI{0}{\percent}.

The mechanism is a systematic over-command. At the instants of the Test~2
gust where the deflection that cancels the excursion is \SI{0.0}{\degree},
\SI{1.4}{\degree} and \SI{4.2}{\degree}, the linear controller commands
\SI{5.1}{\degree}, \SI{8.9}{\degree} and \SI{8.8}{\degree}, against
\SI{0.4}{\degree}, \SI{2.1}{\degree} and \SI{4.9}{\degree} for the LDNet.
The cause is visible in \cref{fig:ablation_mechanism}: the linear model
carries a single flap authority
$\partial C_L/\partial\delta=\SI{0.0253}{\per\degree}$, whereas the
full-order response varies over \SIrange{0.018}{0.047}{\per\degree} with
deflection and state. Underestimating the authority of its own actuator,
the controller asks for more deflection than the section needs, saturates
and rings, and the difference between the lift it predicts and the lift it
obtains reaches \num{0.23}, larger than the excursion it is cancelling.
The deficiency is structural and not a matter of the fitted value: we
refitted the flap gain four ways, one of them matched to the full-order
sensitivity near trim, spanning \SIrange{0.0225}{0.0383}{\per\degree}, and
all four fail, the best of them recovering about one percentage point. No
constant can succeed, because the quantity it approximates varies by a
factor \num{2.6} over the deflection range.

In the comparison of \cref{tab:ablation} the reduced system is the LDNet in
both arms, which gives the LDNet-based controller an exact internal model.
That asymmetry accounts for most of the margin, and we quantify it rather
than assume it small: repeating the comparison with the linear model as the
reduced system reverses the outcome, the linear controller gaining between
\num{28} and \SI{51}{pp} over the LDNet one. Decomposing the four
combinations of internal model and reduced system into a self-consistency
effect and a model effect attributes \SI{48.3}{pp} of the margin to the
controller matching the reduced system and only \SI{8.8}{pp} to the internal
model itself, with the sign of the latter changing between cells.
\Cref{tab:ablation} therefore establishes that the linear internal model
cannot simply be substituted into the loop, but it does not by itself
measure the quality of the two models.

We measure that separately, on held-out full-order trajectories where
neither model is the reference. Each receives the true structural state and
the true deflection history and rolls the same eight-step horizon the
controller uses, and we score the predicted lift coefficient against the
full-order solution. Over \num{1320} probes the LDNet is more accurate by a
factor \num{9.8}, with a \SI{95}{\percent} confidence interval of
\numrange{8.8}{11.0} from a paired bootstrap, and it is the more accurate of
the two at \SI{96.7}{\percent} of the probes.

\Cref{fig:ablation_cfd} stratifies that comparison and locates the region
where the nonlinearity is required. The factor grows from \num{9.2} below
\SI{2}{\degree} of deflection to \num{15.1} above \SI{11}{\degree}, and
from \num{8.8} below \SI{5}{\meter\per\second} of gust amplitude to
\num{20.7} in the \SIrange{25}{35}{\meter\per\second} band: the error of
the linear model grows by a factor \num{2.3} with deflection and \num{2.6}
with gust amplitude, while that of the LDNet does not. The internal model
is therefore responsible for the alleviation where
\cref{sec:introduction} anticipated it would be, at high $r_g$ and large
deflections, where the assumptions of small amplitude and attached flow are
least tenable.

This last comparison measures the prediction the controller consumes over
its horizon, not a controlled trajectory computed with the full-order model
in the loop, which the cost reported in \cref{app:vv} places out of reach.
It also compares against a linear theory without a separation-onset
correction; a dynamic-stall model would be a more demanding term of
comparison, on which the present result does not bear.

\begin{table}[t]
    \centering
    \begin{tabular}{|r|r|r|r|r|r|r|r|}
    \hline
    \rowcolor{bluePoli!40}
    $W_0$ [\si{\meter\per\second}] & $T_g$ [\si{\second}] & $k$ &
    LDNet [\si{\percent}] & linear [\si{\percent}] & difference [\si{pp}] &
    $|\delta|_{\max}$ [\si{\degree}] & pitch ratio \T\B \\
    \hline\hline
    \input{ablation_table_body}
    \hline
    \end{tabular}
    \\[6pt]
    \caption{Internal-model ablation over the eighteen CS-25.341 gust cases:
    peak lift-excursion reduction of \cref{eq:clred} obtained with the
    retained LDNet ($d_s=10$, six-signal input, $L=6$) and with the fitted
    linear unsteady model as internal model of the Model Predictive
    Controller, at equal preview and equal cost function, each arm tuned
    independently with the rule of \cref{subsec:mpc_results}. The last two
    columns report peak deflection and pitch ratio of the linear arm. The
    LDNet column reproduces \cref{tab:cs25_grid} to within \SI{0.05}{pp};
    the linear internal model does not alleviate the gust on any cell. The
    reduced system is the LDNet in both arms, so most of the difference
    reflects the LDNet-based controller having an exact internal model: the
    table shows that the linear model cannot be substituted into the loop,
    not that it is the less accurate of the two, which is measured in
    \cref{fig:ablation_cfd}.}
    \label{tab:ablation}
\end{table}

\begin{figure}[t]
    \centering
    \includegraphics[width=0.6\textwidth]{Images/fig_ablation_mechanism.png}
    \caption{Flap authority $\partial C_L/\partial\delta$ against deflection
    at the trim state, for the retained LDNet ($d_s=10$, six-signal input,
    $L=6$) and for the fitted linear unsteady model, evaluated at frozen
    latent state. The full-order response peaks near
    $\delta=\SI{-4}{\degree}$ and decays on both sides, while a linear model
    carries one constant by construction: underestimating the authority of
    the flap, the Model Predictive Controller that uses it over-commands the
    deflection.}
    \label{fig:ablation_mechanism}
\end{figure}

\begin{figure}[t]
    \centering
    \includegraphics[width=\textwidth]{Images/fig_ablation_cfd.png}
    \caption{Root-mean-square error of the eight-step $C_L$ prediction
    against held-out full-order trajectories, for the retained LDNet
    ($d_s=10$, six-signal input, $L=6$) and the fitted linear unsteady
    model, binned by peak deflection and peak gust amplitude over the
    horizon; \num{1320} probes, both models driven by the true structural
    state and the true deflection history, so that neither is the reference.
    The advantage of the LDNet grows from a factor \num{9.2} to \num{15.1}
    with deflection and from \num{8.8} to \num{20.7} with gust amplitude,
    locating at high $r_g$ and large deflections the region where the
    nonlinearity of the internal model is required.}
    \label{fig:ablation_cfd}
\end{figure}
